% Created 2026-06-20 sam. 14:33
% Intended LaTeX compiler: pdflatex
\documentclass[11pt]{article}

\usepackage[utf8]{inputenc}
\usepackage[T1]{fontenc}
\usepackage{graphicx}
\usepackage{longtable}
\usepackage{wrapfig}
\usepackage{rotating}
\usepackage[normalem]{ulem}
\usepackage{amsmath}
\usepackage{amssymb}
\usepackage{capt-of}
\usepackage{hyperref}
\author{Michaël CHLON}
\date{2026-06-20}
\title{Architecture d'un projet Rust : Package / Crate / Module}
\hypersetup{
 pdfauthor={Michaël CHLON},
 pdftitle={Architecture d'un projet Rust : Package / Crate / Module},
 pdfkeywords={},
 pdfsubject={},
 pdfcreator={},
 pdflang={English}}
\begin{document}

\maketitle
\setcounter{tocdepth}{2}
\tableofcontents

\section{Vue d'ensemble}
\label{sec:org49e60e1}

Rust organise le code sur trois niveaux emboîtés :

\begin{center}
\begin{tabular}{lll}
Niveau & Rôle & Défini par\\
\hline
\textbf{Package} & Unité de distribution / build (Cargo) & \texttt{Cargo.toml}\\
\textbf{Crate} & Unité de compilation (1 arbre de modules) & un fichier racine (.rs)\\
\textbf{Module} & Unité d'organisation + visibilité & mot-clé \texttt{mod}\\
\end{tabular}
\end{center}

Relation d'inclusion :

\begin{verbatim}
Package  ⊇  1+ Crates  ⊇  arbre de Modules  ⊇  items (fn, struct, enum...)
\end{verbatim}
\section{Package}
\label{sec:org5227173}

\begin{itemize}
\item Plus grande unité gérée par \textbf{Cargo}.
\item Décrit par un seul \texttt{Cargo.toml} (nom, version, dépendances, édition).
\item Contient \textbf{au moins une} crate, \textbf{au plus une} crate \emph{library}.
\item Peut contenir \textbf{plusieurs} crates \emph{binary}.
\end{itemize}

Règles de Cargo (convention = racines de crate automatiques) :

\begin{center}
\begin{tabular}{ll}
Chemin & Crate produite\\
\hline
\texttt{src/main.rs} & crate binary (nom du package)\\
\texttt{src/lib.rs} & crate library (nom du package)\\
\texttt{src/bin/*.rs} & crates binary supplémentaires\\
\end{tabular}
\end{center}
\section{Crate}
\label{sec:orgd0134a2}

\begin{itemize}
\item \textbf{Unité de compilation} : le compilateur \texttt{rustc} traite une crate à la fois.
\item Deux types :
\begin{itemize}
\item \textbf{binary} : possède \texttt{fn main()}, produit un exécutable.
\item \textbf{library} : pas de \texttt{main}, produit du code réutilisable (\texttt{.rlib}).
\end{itemize}
\item Chaque crate a un \textbf{module racine} (crate root) :
\begin{itemize}
\item \texttt{main.rs} → module racine de la crate binary
\item \texttt{lib.rs}  → module racine de la crate library
\end{itemize}
\item La crate forme un \textbf{arbre de modules} dont la racine s'appelle \texttt{crate}.
\end{itemize}
\section{Module}
\label{sec:org2016793}

\begin{itemize}
\item Unité d'\textbf{organisation} et de \textbf{contrôle de visibilité} (\texttt{pub}).
\item Déclaré avec \texttt{mod}. Par défaut tout est \textbf{privé} (visible seulement dans le module et ses descendants).
\item Trois façons de définir le contenu d'un module :
\end{itemize}

\begin{verbatim}
// 1. inline
mod jardin {
    pub mod legumes { pub fn planter() {} }
}

// 2. fichier frère : src/jardin.rs
mod jardin;

// 3. dossier : src/jardin/mod.rs  (ou src/jardin.rs + src/jardin/legumes.rs)
mod jardin;
\end{verbatim}
\subsection{Chemins (paths)}
\label{sec:orga1476c2}

\begin{itemize}
\item \textbf{absolu} : part de \texttt{crate} → \texttt{crate::jardin::legumes::planter()}
\item \textbf{relatif} : \texttt{self}, \texttt{super} (module parent), ou nom direct.
\item \texttt{use} crée un raccourci (import) dans le scope courant.
\end{itemize}

\begin{verbatim}
use crate::jardin::legumes::planter;
planter();
\end{verbatim}
\subsection{Visibilité}
\label{sec:org8ce9101}

\begin{center}
\begin{tabular}{ll}
Mot-clé & Portée\\
\hline
(rien) & privé au module\\
\texttt{pub} & public\\
\texttt{pub(crate)} & public dans toute la crate seulement\\
\texttt{pub(super)} & public pour le module parent\\
\end{tabular}
\end{center}
\section{Schéma explicatif}
\label{sec:org0f16377}

\begin{verbatim}
                         ┌──────────────────────────────────────────┐
                         │              PACKAGE                     │
                         │            (Cargo.toml)                  │
                         │   nom · version · édition · deps         │
                         └──────────────────────────────────────────┘
                                          │ contient
                  ┌───────────────────────┴────────────────────────┐
                  ▼                                                 ▼
        ┌───────────────────┐                          ┌───────────────────────┐
        │  CRATE library    │                          │   CRATE binary        │
        │   src/lib.rs      │                          │    src/main.rs        │
        │  (crate root)     │                          │   fn main() {…}       │
        └───────────────────┘                          └───────────────────────┘
                  │ arbre de modules
                  ▼
        crate (racine)
        ├── mod reseau            (src/reseau.rs)
        │   ├── pub fn connecter()
        │   └── mod client        (src/reseau/client.rs)
        │       └── pub fn envoyer()
        └── mod jardin            (src/jardin/mod.rs)
            └── pub mod legumes
                └── pub fn planter()

   Accès depuis main.rs :
       use ma_lib::reseau::connecter;   // crate externe (la library du package)
       crate::jardin::legumes::planter(); // chemin absolu interne
\end{verbatim}
\section{Arborescence type d'un package}
\label{sec:org105b0fa}

\begin{verbatim}
mon_package/
├── Cargo.toml          ← définit le PACKAGE
├── src/
│   ├── lib.rs          ← racine de la CRATE library
│   ├── main.rs         ← racine d'une CRATE binary
│   ├── reseau.rs       ← MODULE reseau
│   ├── reseau/
│   │   └── client.rs   ← MODULE reseau::client
│   └── bin/
│       └── outil.rs    ← CRATE binary supplémentaire
└── tests/              ← crates d'intégration (1 crate par fichier)
\end{verbatim}
\section{À retenir}
\label{sec:orgfdd8a76}

\begin{itemize}
\item \texttt{Cargo.toml} = 1 package.
\item 1 package = 0..1 library + 0..N binaries.
\item 1 crate = 1 compilation = 1 arbre de modules avec racine \texttt{crate}.
\item \texttt{mod} découpe, \texttt{pub} expose, \texttt{use} raccourcit, les paths localisent.
\end{itemize}
\end{document}
